\section{Problem}

\begin{definicja}
Automorfizmem wielomianowym płaszczyzny $\Q^2$ nazywamy odwzorowanie
$F=(F_1,F_2)\colon\Q^2\to\Q^2$, którego składowe są wielomianami o współczynnikach
wymiernych i które ma odwrotność $G=(G_1,G_2)$ tej samej postaci. Stopniem $F$ nazywamy
$\deg F=\max(\deg F_1,\deg F_2)$, gdzie $\deg$ oznacza stopień całkowity.
\end{definicja}

Dane są parami różne punkty $p_1,\dots,p_m\in\Q^2$ oraz parami różne punkty
$q_1,\dots,q_m\in\Q^2$. Pytamy o najmniejszy stopień automorfizmu $F$ spełniającego
$F(p_i)=q_i$ dla $i=1,\dots,m$. Wiadomo, że stopień $(m-1)^2$ zawsze wystarcza.
Celem eksperymentu jest zbadanie, czy dla konkretnej bijekcji pięciu punktów można
zejść niżej, i -- o ile się da -- uzasadnienie, że niżej zejść się nie da.
